\documentclass[10pt, letterpaper]{article}

\usepackage{../../../../template/template}

% ------------------------
% Impostazioni documento
% ------------------------

\setDocTitle{Verbale esterno 01/04/2026}
\setDocVersion{-}
\setDocData{01/04/2026}
\setDocIndex{ve\_2026\_04\_01}
\setDocState{2}
\setDocRecipient{2}

\begin{document}

\makeTitlePage

\newpage

\tableofcontents

\newpage

% -------------------------------------------------
% Informazioni riunione
% -------------------------------------------------
\makeMeetingInfoTable{01/04/2026}{16:00}{17:30}{via \textit{Microsoft Teams$_G$}}

% ------------------------
% Partecipanti
% ------------------------
\addParticipant{Filippo}{Venzo}{Programmatore}{P}
\addParticipant{Valeria}{Baleanu}{Programmatore}{P}
\addParticipant{Luca}{Granziero}{Amministratore}{P}
\addParticipant{Christian}{Libralato}{Responsabile}{P}
\addParticipant{Francesco}{Pasqual}{Programmatore}{P}
\addParticipant{Leonardo}{Pellizzon}{Programmatore}{P}
\addParticipant{Giuseppe}{De Fina}{Verificatore}{P}
\makeMeetingParticipantsTable

% -------------------------------------------------
% Ordine del giorno
% -------------------------------------------------
\section{Ordine del giorno}
\begin{itemize}
      \item Revisione bozzetti con il team UX/UI;
      \item pianificazione delle Sessioni di Demo e Presentazione;
      \item implementazione delle Analytics e dei Suggerimenti;
      \item rinegoziazione e Definizione dei Requisiti di Dashboard e Interazione;
      \item preparazione al Collaudo e Deliverables di Progetto;
      \item visualizzazione interfaccia grafica;
      \item tipologia dei Grafici Analytics;
      \item simulazione e Test dei Sensori di Caduta.
\end{itemize}

% -------------------------------------------------
% Approfondimento
% -------------------------------------------------
\section{Approfondimento}

\subsection*{Revisione bozzetti con il team UX/UI}
In data 30/03/2026 è stata effettuata una riunione con la Proponente che ha previsto la partecipazione del team UX/UI al fine di revisionare, commentare 
e fornire suggerimenti rispetto ai bozzetti grafici realizzati da \teamName. Il Responsabile ha esposto e descritto le varie sezioni dei bozzetti, che corrispondevano alle 
principali sezioni dell'applicativo, motivando le scelte stilistiche e strutturali.
In riferimento a ciò il team UX/UI ha fornito molti suggerimenti e fatto emergere alcune criticità funzionali e di design, le principali sono state:
\begin{itemize}
      \item la necessità di rendere chiaro all'interno della dashboard quali widget siano fissi e quali no;
      \item dare importanza agli allarmi nella dashboard e prestare attenzione all'ordine con cui vengono mostrati, un ordinamento per priorità si adatta maggiormente al caso d'uso;
      \item le informazioni mostrate dalla mappa nella sezione dispositivi sono ridondanti e tolgono spazio all'elenco dispositivi;
      \item valutare altri posizionamenti rispetto alla barra laterale per i suggerimenti nella sezione analytics in modo da fornire più importanza;
      \item valutare la possibilità di creare un unico selettore di impianto che mantenga coerenza tra le varie sezioni dell'applicativo piuttosto che avere selettori indipendenti in pagine diverse;
      \item modernizzare l'interfaccia e rivalutare i colori per ridurre la presenza, elevata, del bianco.
\end{itemize}
Il team si è prefissato l'obiettivo di applicare per quanto possibile i consigli emersi, pur considerando che la codifica si trova ad uno stato intermedio di completamento e il tempo rimasto è ridotto a causa 
delle scadenze imminenti.

\subsection*{Pianificazione delle Sessioni di Demo e Presentazione}
Sono state proposte due sessioni per il collaudo: una tecnica di un'ora e mezza in presenza l'8 aprile e una demo di mezz'ora in remoto il 9 aprile,
coinvolgendo sia il team tecnico che il reparto marketing, con la possibilità di mostrare i risultati e spiegare il funzionamento della dashboard. 
La Proponente si è inoltre resa disponibile a fornire un server qualora necessario, richiedendo specifiche tecniche per la configurazione 
e sottolineando l'importanza di avere tutto pronto per l'8 aprile.

\subsection*{Implementazione delle Analytics e dei Suggerimenti}
È stata esposta l'implementazione delle analytics e dei suggerimenti tramite il modello Groq, spiegando la logica di caching, la periodicità di aggiornamento.
Il modello Groq viene utilizzato per generare suggerimenti gratuiti, con limiti di 30 richieste al minuto e 14.000 al giorno, integrando i suggerimenti direttamente nelle analytics inviate dal backend al frontend.
I suggerimenti sono inclusi nelle analytics, la generazione di entrambi avviene on demand tramite trigger sull'accesso alla sezione Analytics, con caching gestito da \textit{nginx$_G$} per utente e Plant ID,
e scadenza della cache ogni 24 ore; non è previsto un refresh forzato, ma è stata caldamente consigliata un'opzione di refresh manuale. 

\subsection*{Rinegoziazione e Definizione dei Requisiti di Dashboard e Interazione}
Si è discussa la necessità di rinegoziare alcuni requisiti opzionali, proponendo di rendere la dashboard fissa e non personalizzabile e rimuovere la mappa degli allarmi.

È stato poi concordato di mostrare i suggerimenti nei grafici tramite un box laterale, evitando la sidebar per valorizzare i suggerimenti e facilitare la comprensione del riferimento metrico. Infine è stata
sottolineata l'importanza di includere la location nell'allarme per permettere all'operatore di sapere dove intervenire, compensando la rimozione della mappa degli allarmi, oltre che dare importanza agli allarmi 
ad alta priorità nella dashboard.

\subsection*{Preparazione al Collaudo e Deliverables di Progetto}
Sono state illustrate le modalità di collaudo, i deliverables richiesti e le aspettative per la presentazione finale, specificando la necessità di mostrare l'avanzamento, la demo, l'architettura (modello C4), la timeline e una retrospettiva, con indicazioni sulla gestione dei bozzetti e dei manuali utente.
È stato inoltre chiarito che i bozzetti possono essere presentati come versione preliminare, senza necessità di aggiornarli completamente, ma è auspicabile mantenere coerenza tra bozzetti e implementazione finale.

\subsection*{Visualizzazione interfaccia grafica}
È stato presentato lo stato dell'applicativo, con particolare riferimento alla gestione degli utenti e dei reparti, illustrando le soluzioni adottate per la visualizzazione e la creazione tramite form,
con discussione sulla coerenza grafica, l'uso dei colori e l'integrazione con l'autenticazione interna e verso Vimar.

\subsection*{Tipologia dei Grafici Analytics}
Si è poi discussa la scelta tra grafici a barre e a punti per le diverse analytics, concordando sulla visualizzazione in base alla continuità dei dati.
È stato concordato di inserire un selettore per appartamento e di indicare chiaramente il periodo di riferimento dei dati, come gli ultimi 30 giorni, per facilitare la comprensione da parte dell'utente.

\subsection*{Simulazione e Test dei Sensori di Caduta}
Si è infine configurata la simulazione e il test dei sensori di caduta, inviando l'account amministratore per l'aggiunta a un impianto di test, con la possibilità di simulare eventi tramite CURL e verificare la corretta ricezione degli allarmi.

% -------------------------------------------------
% Decisioni
% -------------------------------------------------
\addDecision{Aggiunta della funzionalità di refresh manuale per le Analytics}
\addDecision{Spostamento dei suggerimenti in un box laterale associato al grafico con scorrimento verticale}
\addDecision{Riduzione requisiti opzionali, con mantenimento di dashboard statica e rimozione mappa allarmi}
\makeDecisionTable

% -------------------------------------------------
% Azioni
% -------------------------------------------------
\addTodo{\noOne}{Redigere verbale esterno}{C. Libralato}{05/04/2026}
\addTodo{\#234}{Preparare presentazione per il collaudo.}{\teamName}{08/04/2026}
\addTodo{\noOne}{Decidere la piattaforma di deployment.}{\teamName}{08/04/2026}

\makeTodoTable

\end{document}